\chapter{Requirement Analysis}

This chapter records the requirements that the system was built to satisfy,
separated into functional behaviour, quality attributes, and the environment
assumed by the implementation.

\section{Functional Requirements}

\tabref{tab:funcreq} lists the functional requirements, each with the
identifier used elsewhere in the report.

\begin{table}[H]
\centering
\caption{Functional requirements}
\label{tab:funcreq}
\small
\begin{tabular}{@{}p{1.4cm}p{8.4cm}p{2.6cm}@{}}
\toprule
\textbf{ID} & \textbf{Requirement} & \textbf{Priority} \\
\midrule
FR-1 & A device or controller shall enroll using a single-use code and thereafter be identified by an Ed25519 public key. & Essential \\
FR-2 & The signaling server shall authenticate every connection by a challenge--response signature before accepting any other message. & Essential \\
FR-3 & A host shall register its availability and be addressable by a nine-digit connect identifier. & Essential \\
FR-4 & A controller shall initiate a session by connect identifier and the host shall be able to accept or reject it. & Essential \\
FR-5 & The peers shall exchange SDP offers, answers and ICE candidates through the signaling server and establish a direct peer connection. & Essential \\
FR-6 & The host shall capture the primary display and transmit it as an encoded video stream. & Essential \\
FR-7 & The controller shall transmit keyboard and pointer events, which the host shall inject into the local input queue. & Essential \\
FR-8 & The native host shall operate with no window, no user gesture and no source picker. & Essential \\
FR-9 & The native host shall accept unattended sessions automatically on presentation of a configured password. & Essential \\
FR-10 & The native host shall install as a Windows service, start at boot and survive session switching. & Essential \\
FR-11 & The signaling server shall optionally persist enrolled identities and a session audit log to PostgreSQL. & Desirable \\
FR-12 & The signaling server shall optionally terminate TLS and serve \texttt{wss://}. & Desirable \\
FR-13 & An administrative utility shall list enrolled devices with their connect identifiers. & Desirable \\
\bottomrule
\end{tabular}
\end{table}

\section{Non-Functional Requirements}

\subsection{Performance}

Interactive remote control is unusable if the delay between an input event and
the corresponding visual change is perceptible as lag. The target adopted was an
end-to-end glass-to-glass latency below 150\,ms on a local network and below
250\,ms over a typical broadband path, with a sustained frame rate of at least
25\,frames per second at 1920$\times$1080. Bandwidth consumption was to remain
within the capability of an ordinary residential uplink for typical desktop
content.

\subsection{Security}

The signaling server must not be able to decrypt session media. Endpoint
identity must rest on private keys that never leave the machine holding them.
Enrollment must not be open: possession of a valid one-time code is required
before a key is accepted. Unattended acceptance must be gated by a secret that
is verified by the host rather than by the server.

\subsection{Reliability and Availability}

The service-mode host must return to service automatically after a reboot,
after a user logs off, and after the console session is switched. Failure of the
signaling server must not terminate sessions that have already established a
direct peer connection, since the server is not on the media path.

\subsection{Maintainability}

The system is organised as a Cargo workspace with separate crates for the
signaling server, the native host and the service supervisor, so that each can
be built and tested in isolation. Protocol messages are JSON, which keeps them
inspectable during development.

\section{Hardware and Software Environment}

\begin{table}[H]
\centering
\caption{Development and test environment}
\label{tab:env}
\small
\begin{tabular}{@{}p{4.2cm}p{9cm}@{}}
\toprule
\textbf{Component} & \textbf{Specification} \\
\midrule
Host machine        & Windows 11 Pro (build 26200), x86-64, 16\,GB RAM \\
Controller machine  & Windows 11, x86-64, 8\,GB RAM \\
Language toolchain  & Rust (2021 edition), stable channel \\
Async runtime       & Tokio \\
WebRTC stack        & \texttt{webrtc-rs} 0.11, with a patched \texttt{webrtc-dtls} \\
Video encoder       & OpenH264 via the \texttt{openh264} crate \\
Screen capture      & \texttt{scrap} (Desktop Duplication) \\
Input injection     & \texttt{enigo}, and the Win32 \texttt{SendInput} path \\
Signaling transport & \texttt{tokio-tungstenite} over WebSocket / WSS \\
Cryptography        & \texttt{ed25519-dalek} 2.x \\
TLS                 & \texttt{tokio-rustls} with the \texttt{ring} provider \\
Persistence         & PostgreSQL via \texttt{sqlx} (optional) \\
GUI applications    & Tauri with a WebView2 front end \\
Network             & 1\,Gbps switched LAN; consumer broadband with CGNAT for wide-area tests \\
\bottomrule
\end{tabular}
\end{table}

\section{Feasibility Study}

\textbf{Technical.} Every component has a viable Rust implementation, which was
confirmed by prototyping the riskiest element first --- silent framebuffer
capture --- before committing to the architecture. The one genuine technical
obstacle encountered was a DTLS handshake failure against current Chrome
versions, resolved as described in Section~5.6.

\textbf{Economic.} All dependencies are open source. The only recurring cost is
a modest virtual private server for signaling and an optional TURN relay, and
because media is peer-to-peer the relay is exercised only in the minority of
network configurations that defeat direct connectivity.

\textbf{Operational.} The host installs as a service with a single command and
requires no configuration beyond an enrollment code and an unattended password,
which is comparable to the deployment effort of the commercial alternatives.
